\documentclass[12pt,a4paper]{article}

% Packages (same font and layout as the CV, main.tex)
\usepackage[utf8]{inputenc}
\usepackage[T1]{fontenc}
\usepackage{newpxtext,newpxmath} % Use Palatino font
\usepackage[margin=1in]{geometry}
\usepackage{titlesec}
\usepackage{hyperref}
\usepackage{xcolor}

% Formatting
\pagestyle{empty} % No page numbers
\setlength{\parindent}{0pt} % No indentation
\setlength{\parskip}{6pt} % Space between paragraphs
\clubpenalty=10000 \widowpenalty=10000 % No single lines stranded at a page break

% Custom Section Formatting
\titleformat{\section}{\large\bfseries}{}{0em}{}[\titlerule]
\titlespacing{\section}{0pt}{12pt}{6pt}

% Hyperlink Setup
\hypersetup{
    colorlinks=true,
    linkcolor=blue,
    filecolor=magenta,
    urlcolor=blue,
}

% Red placeholder for information still to be filled in (remove all uses before submitting)
\newcommand{\todo}[1]{\textcolor{red}{[#1]}}

\begin{document}

%----------------------------------------------------------------------------------------
%	HEADER
%----------------------------------------------------------------------------------------
\begin{center}
    {\LARGE \textbf{Supporting Statement}} \\ \vspace{5pt}
    {\large Haotian Zhang} \\ \vspace{5pt}
    Postdoctoral Research Assistant in Quantitative Finance (Vacancy 188851) \\ \vspace{3pt}
    E-Mail: \url{htzhang@link.cuhk.edu.hk}
\end{center}

%----------------------------------------------------------------------------------------
%	OVERVIEW
%----------------------------------------------------------------------------------------
I am a financial econometrician working with ultra-high-frequency financial data. 
Transaction prices differ from efficient prices because of market frictions
such as microstructure noise and price staleness.
My research develops estimators and inference procedures for 
these frictions within the trading day 
and that remain valid in the episodes where 
the market is turbulent and illiquid, such as during 
macroeconomic announcements, closing auctions, and flash crashes. 
Frictions determine what it costs to trade 
and how quickly prices incorporate information,
so measuring them accurately is a first step toward testing return predictability
at short horizons and evaluating trading strategies net of costs. 
Therefore, my research is highly relevant to the financial economics and financial data science research 
at the Oxford-Man Institute (OMI).

Below I explains how I meet each of the \textit{essential selection criteria}
at the OMI.

\textbf{1. Relevant PhD, with a focus on financial data science and market microstructure.}
I am completing a Ph.D. in Economics at The Chinese University of Hong Kong (2021--present),
supervised by Prof.~Z. Merrick Li, with a thesis in financial econometrics. 
My thesis is titled \textit{``Beyond Itô Semimartingales:
Estimation and Inference for Market Frictions''}. 
The abstract of the thsis was approved by my Ph.D. committee 
and submitted to the Chinese University of Hong Kong in October 2026. 
Thus, I am expected to graduate before August 2027. 
The thesis is a comprehensive study of market microstructure and financial data science 
including the three working papers listed in my CV. 

My work is data-driven throughout: 
each of my papers starts from a new feature of
tick-by-tick transaction data and ends with a new estimator (or inference procedure) 
and rich empirical applications. 
For example, 
in the paper \textit{Estimation and Inference for Spot Variance of Microstructure Noise}, 
we estimate the intraday pattern of the spot noise variance 
from SPY transactions over 2016--2025 
and detect the scheduled jumps at key macroeconomic news releases 
and closing-auction deadlines. 
Before the Ph.D., I completed a B.Sc.\ in Mathematical Economics and Statistics
at the University of Birmingham, where I received the Economic Final Year Prize
(top two in the cohort); and a Bachelor in Economics (with minor in finance) 
at the Southwestern University of Finance and Economics, 
where I received the National Scholarship in 2019. 

\textbf{2. Specialist knowledge of the discipline.}
My specialist field is high-frequency financial econometrics and market microstructure.
The papers listed in my CV combine limit theory for semimartingales,
including stable central limit theorems and strong Gaussian approximations,
with hands-on work on ultra-high-frequency transaction data.
This combination fits the OMI's data-driven research into how markets work
and its work on market microstructure and high-frequency trading.

\textbf{3. Ability to manage my own research and associated activities.}
I am the sole author of ``Uniform Inference for Intraday Market Frictions'',
for which I carried out every stage myself, from the research question and the asymptotic theory
to the simulations and the empirical application.
Alongside it, I work on two coauthored projects with 
Prof. Merrick Li at The Chinese University of Hong Kong, 
Prof. Federico M. Bandi at Johns Hopkins University, 
and Prof. Roberto Ren\`o at ESSEC Business School,
which requires planning work and meeting deadlines 
across three institutions and time zones.
These experiences have given me a strong ability to conduct independent research, 
and collarate with colleagues from different institutions.

\textbf{4. Experience of contributing to publications and presentations.} 
The paper \textit{Estimation and Inference for Spot Variance of Microstructure Noise}
is now 
under review at the \textit{Journal of the American Statistical Association}.
I presented my work at the 17th Annual Meeting of the Society for Financial Econometrics
(SoFiE; Paris, June 2025), 
the 2025 World Congress of the Econometric Society (Seoul, August 2025),
and the 
First Research Colloquium in Financial Econometrics (ESSEC Business School, May 2026). 
I also serve as a referee for the \textit{Journal of Financial Econometrics}, 
which gives me a deep understanding of the research process 
and the expectations of relevant journals.


\textbf{5. Communication skills.}
I have written for publication as a coauthor of a published article
and of a paper now under review, and as the sole author of a working paper.
I have presented to specialist audiences at the three international meetings listed above.
At SoFiE2025 and the ESWC2025 I presented the joint paper with my supervisor
on behalf of both authors, which is the role I would take
when representing the research group at meetings.
Refereeing has also given me practice in assessing the work of others
and explaining that assessment clearly.

\textbf{6. Respect, courtesy, and an inclusive environment.}
I have studied and worked in four academic communities,
in Chengdu, Birmingham, Hong Kong, and Paris,
and my coauthors work in three countries and are at different career stages.
Working across these settings depends on treating colleagues with respect and courtesy,
on listening before disagreeing, and on giving credit fairly.
I would bring the same conduct to the OMI, and I am committed to keeping
the Department a positive, supportive, and inclusive environment for all.

\end{document}
